\documentclass[11pt,a4paper]{article}
\usepackage[utf8]{inputenc}
\usepackage[T1]{fontenc}
\usepackage{amsfonts}
\usepackage[french]{babel}
\frenchbsetup{StandardLists=true} % à inclure si on utilise \usepackage[francais]{babel}
\usepackage{enumitem}
\usepackage{amssymb}
\usepackage{amsmath}
\usepackage[left=2cm,right=2cm,top=2cm,bottom=2cm]{geometry}
\usepackage{multirow}
\usepackage[dvipsnames]{xcolor}

\usepackage[T1]{fontenc}
\usepackage{fouriernc}
\usepackage{graphicx}

\usepackage{setspace}
\singlespacing

\usepackage{pstricks}
\usepackage{xkeyval}
\usepackage{pst-labo}

\usepackage{tcolorbox}
\usepackage{multicol}
\usepackage{caption}
\usepackage{fancyhdr}
\pagestyle{fancy}
\renewcommand\headrulewidth{1pt}
\fancyhead[L]{Lycée Denis Diderot }
\fancyhead[R]{Classe de Première Bac Pro}
\setcounter{page}{2}\fancyfoot[C]{page \thepage}
\fancyfoot[R]{Année 2025-2026}
\fancyfoot[L]{Alexis RUIZ}
\usepackage{multirow,tabularx,booktabs}
\newcommand{\cadreblanc}[1]{%
\medskip\framebox[\linewidth][c]{%
\rule{0mm}{#1mm}}\par
\medskip}

\usepackage{awesomebox}
\setlength{\aweboxvskip}{0mm}
\usepackage{pifont}
\usepackage{chemmacros}
\chemsetup{modules=all}
\setlength{\columnseprule}{.25mm}
\setlength{\parindent}{0pt}

\renewcommand{\arraystretch}{2}
\newcommand*\acocher{\quitvmode{\fboxsep0pt \fboxrule0.8pt \fbox{\vrule height1.8ex depth.1ex width0pt \vrule height0pt depth0pt width1.9ex }}}

\newcommand{\code}{\awesomebox[black]{2pt}{

\includegraphics[scale=0.15]{Chapitre 1 - Energie et puissance électrique/images/frame.png}
}{black}}

\newcommand{\moteur}{\awesomebox[black]{2pt}{

\includegraphics[scale=1]{images/moteur.jpg} 
}{black}}

\newcommand{\defconclusion}{\awesomebox[black]{2pt}{\rotatebox{90}{\Large{\textbf{Conclusion}}}}{black}}

\frenchbsetup{StandardLists=true}



\begin{document}


\hrule\


\begin{center}
 \LARGE{Activité 1 - Revoir l'intensité électrique}
\end{center}


\normalsize

\hrule\
\\[0,2cm]



Une analogie intéressante pour comprendre de façon simple les notions d'intensité et de tension peut être faite avec l'écoulement d'une rivière. Celle-ci s'écoule d'amont vers l'aval.  Le dénivelé, différence d'altitude entre le point haut et le point bas de la rivière, peut être assimilé à la tension, le débit à l'intensité du courant et la taille à la section d'un câble électrique.


\newtcolorbox{mybox}{colback=white,
colframe=black!5!black}

\vfill 

\begin{enumerate}
    \item Mesure de l'intensité :

\cadreblanc{30}



Exercice 1 : Convertir (avec 1A = 1000 mA)  ~~ 


\vspace{1mm}

\begin{multicols}{4}
$1848$ mA = \dotfill A 
\vspace{5mm}

$77$ mA = \dotfill A 
\vspace{5mm}


$1,085$ A = \dotfill mA
\vspace{5mm}


$0,021$ A = \dotfill mA
\vspace{5mm}
    
\end{multicols}

\cadreblanc{25}
    \vfill

\item Lois de l'intensité : \\[-2mm]
    \vspace{-2mm}
    \begin{itemize}
        \item Dans un circuit en série, l'intensité est la même en tout point du circuit. 
        \item L'intensité qui sort du générateur est identique à celle qui y retourne, conformément à la loi de Kirchhoff.
        \item Dans un circuit en dérivations, l'intensité dans la branche principale est égale à la somme des intensités dans les branches dérivées. Il s'agit de la loi des noeuds. 
    \end{itemize}
    \vspace{-3mm}
    
    \begin{center}
    
        \begin{minipage}{7cm}
        \cadreblanc{10}
        \end{minipage}
        
    \end{center}

    
    Exercice 2 :
    Dans le circuit ci dessous \\
    \vspace{-5mm}
\noindent%
\begin{minipage}{.65\textwidth}%
 On donne $\mathrm{I_2 }$ = 0,7 A et $\mathrm{I_3 }$ = 0,3 A.
 
En appliquant la loi des noeuds, calculer les
intensités manquantes. 

\renewcommand{\arraystretch}
{3}
\begin{tabular}{|c||p{1.cm}|p{1.cm}|p{1.cm}|p{1.cm}|p{1.cm}|p{1.cm}|} \hline
Intensité  &  $I_1$& $I_2$ &$I_3$  & $I_4$ & $I_5$ & $I_6$ \\ \hline
Valeur (A)& & & & & &\\ \hline \end{tabular}



\end{minipage}%
\hfill
\begin{minipage}{.3\textwidth}%
 \begin{center}
     

\includegraphics[scale=0.4]{Chapitre 1 - Energie et puissance électrique/images/ex2.png} \end{center}
\end{minipage}%

\vfill 




          

 



    




\end{enumerate}





\end{document}